\documentclass[11pt]{book}
\usepackage[margin=1in]{geometry}
%\usepackage[utf8]{inputenc}
\usepackage{amsmath}
%\usepackage[T1]{fontenc}
\usepackage{float}
\usepackage{amsfonts}
\usepackage{amssymb}
\usepackage{graphicx}
\usepackage[french]{babel}
\author{Christopher Blier-Wong}
\title{Étude examen prédoctoral général}
%\date{Version du \today}
\date{\vfill\includegraphics[width=0.4\textwidth]{UL_actu_3l}}
\usepackage{hyperref}
\hypersetup{colorlinks}
\usepackage{luatexbase}
\usepackage{amsthm}
\usepackage{tikz}
\usepackage[figurename=Illustration]{caption}
\usepackage[tablename=Tableau]{caption}
\usepackage{enumitem}
\usepackage{pbox}
\usepackage[most]{tcolorbox}
\tcbuselibrary{theorems}
\tcbuselibrary{skins}
\usetikzlibrary{shadings}
\usepackage{multicol}
\usepackage{cancel}
\usepackage{microtype}
\usepackage{vwcol} 
\usepackage{bbm}
\usepackage{actuarialsymbol}

\definecolor{dlblue}{RGB}{43,252,226}
\definecolor{dlblack}{HTML}{282727}
\definecolor{dlbeige}{HTML}{FCFCF5}
\definecolor{dlred}{RGB}{255,107,107}
\definecolor{dlyellow}{RGB}{255, 230, 109}

\usepackage{hyperref}
\hypersetup{
	colorlinks = true,
	linkcolor = dlred,
	urlcolor  = dlred,
	citecolor = dlred,
	anchorcolor = dlred
}

\makeatletter
\@addtoreset{chapter}{part}
\makeatother  

\usepackage{fontspec}
\usepackage{titling}

\setmainfont[Path=Muli/, BoldFont=Muli-Bold.ttf,
ItalicFont=Muli-Italic.ttf,
BoldItalicFont=Muli-BoldItalic.ttf]{Muli-Regular.ttf}%

\usepackage{setspace}
\usepackage{pgfplots}
\usepackage{titlesec, blindtext, color}
\definecolor{gray75}{gray}{0.75}
\newcommand{\hsp}{\hspace{20pt}}
\newcommand{\hspsection}{\hspace{10pt}}
\newcommand{\hspsubsection}{\hspace{5pt}}
\titleformat{\chapter}[hang]{\Huge\bfseries}{\thechapter\hsp\textcolor{gray75}{|}\hsp}{0pt}{\Huge\bfseries}
\titleformat{\section}[hang]{\Large\bfseries}{\thesection\hspsection\textcolor{gray75}{|}\hspsection}{0pt}{\Large\bfseries}
\titleformat{\subsection}[hang]{\large\bfseries}{\thesubsection\hspsubsection\textcolor{gray75}{|}\hspsubsection}{0pt}{\large\bfseries}
\usepackage{polynom}
%\input{longdiv.tex}
\usetikzlibrary{spy}
\usepgfplotslibrary{groupplots, colormaps}
\usetikzlibrary{matrix}

\pgfplotsset{
	colormap={dlredblue}{
		rgb255={255,107,107}
		rgb255=(43,252,226)
	},
}
\pgfplotsset{soldot/.style={color=dlblue,only marks,mark=*}} 
\pgfplotsset{holdot/.style={color=dlblue,fill=white,only marks,mark=*}}
\pgfplotsset{
	integral segments/.code={\pgfmathsetmacro\integralsegments{#1}},
	integral segments=3,
	integral/.style args={#1:#2}{
		ybar interval,
		domain=#1+((#2-#1)/\integralsegments)/2:#2+((#2-#1)/\integralsegments)/2,
		samples=\integralsegments+1,
		x filter/.code=\pgfmathparse{\pgfmathresult-((#2-#1)/\integralsegments)/2}
	}
}
\usepackage{fancyhdr}
\pagestyle{fancyplain}
\fancyhf{}
%\setlength{\headheight}{52pt}
\fancyhead[LO]{\leftmark}
\fancyhead[RE]{\rightmark}
\fancyhead[LE,RO]{\thepage}
%\fancyhead[REF]{Version du \today}
%\rfoot{Version du \today}
\usepackage{cleveref}
\tikzstyle{arrow} = [thick,->,>=stealth]
\usetikzlibrary{arrows.meta}
\usetikzlibrary{arrows,decorations.markings}
\crefformat{section}{\S#2#1#3}
\crefformat{subsection}{\S#2#1#3}
\crefformat{subsubsection}{\S#2#1#3}

\usetikzlibrary{arrows, shapes}
\usepgfplotslibrary{polar}
\usetikzlibrary{calc,patterns,angles,quotes}

\newcommand{\laplace}[1]{\mathcal{L}\left\{#1\right\}}
\newcommand{\inverselaplace}[1]{\mathcal{L}^{-1}\left\{#1\right\}}
\newcommand{\exercicemath}[1]{\displaystyle \boldsymbol{#1}}
\newcounter{mytheorem}[section]
\newcounter{myexample}[section]
\newcounter{mytable}[chapter]

\def\themytheorem{\thesection.\arabic{mytheorem} }
\def\themyexample{\thesection.\arabic{myexample} }
\def\themytable{\thechapter.\arabic{mytable} }

\tcbset{
	defstyle/.style={fonttitle=\bfseries\upshape,
		arc=2mm, colback=dlblue!10,colframe=dlblue, coltitle=dlblack, breakable,pad at break=5mm,},
	propstyle/.style={fonttitle=\bfseries\upshape,
		arc=2mm, colback=dlblue!10,colframe=dlblue, coltitle=dlblack, breakable,pad at break=5mm,},
	theostyle/.style={fonttitle=\bfseries\upshape, fontupper=\slshape,
		arc=2mm, colback=dlred!10,colframe=dlred, coltitle=dlblack, breakable, pad at break=5mm,},
	notestyle/.style={fonttitle=\bfseries\upshape,
		arc=2mm, colback=dlyellow!10,colframe=dlyellow, coltitle=dlblack, breakable},
	astucestyle/.style={fonttitle=\bfseries\upshape, fontupper=\slshape,
		arc=2mm, colback=dlbeige,colframe=dlblack},
	examplestyle/.style={skin=standard jigsaw,fonttitle=\bfseries\upshape, 
		arc=2mm, opacityback=0, breakable,pad at break=1mm},
	exampleimagestyle/.style={skin=standard jigsaw,fonttitle=\bfseries\upshape, 
		arc=2mm, opacityback=0, breakable,pad at break=1mm, },
	%	examplestyle/.style={skin=standard jigsaw,fonttitle=\bfseries\upshape, fontupper = \normalsize, 
	%		arc=2mm, colback = white, breakable, opacitytext=0.0,tikz={opacity=0}, coltext=white, 
	%		pad at break=1mm,},
	tablestyle/.style={skin=standard jigsaw,fonttitle=\bfseries\upshape,
		arc=2mm, opacityback=0},
}

\tcbmaketheorem{definition}{Définition}{defstyle}{mytheorem}{def}
\tcbmaketheorem{proposition}{Proposition}{propstyle}{mytheorem}{prop}
\tcbmaketheorem{propriete}{Propriété}{propstyle}{mytheorem}{propriete}
\tcbmaketheorem{propositionpreuve}{Proposition}{propstyle}{mytheorem}{prop}
\tcbmaketheorem{theoreme}{Théorème}{theostyle}{mytheorem}{thm}
\tcbmaketheorem{theoremepreuve}{Théorème}{theostyle}{mytheorem}{thm}
\tcbmaketheorem{corrolaire}{Corollaire}{theostyle}{mytheorem}{cor}
\tcbmaketheorem{note}{Note}{notestyle}{mytheorem}{note}
\tcbmaketheorem{remarque}{Remarque}{notestyle}{mytheorem}{rem}
\tcbmaketheorem{astuce}{Astuce}{astucestyle}{mytheorem}{astuce}
\tcbmaketheorem{example}{Exemple}{examplestyle}{myexample}{ex}
\tcbmaketheorem{examplelower}{Exemple}{examplestyle}{myexample}{ex}
\tcbmaketheorem{exampleimage}{Exemple}{exampleimagestyle}{myexample}{ex}
\tcbmaketheorem{exampleretirer}{Exemple}{exampleretirerstyle}{myexample}{ex}
\tcbmaketheorem{tablebox}{Tableau}{tablestyle}{mytable}{tab}

\usepackage{pgfplots}
\pgfplotsset{%
	every tick label/.append style = {font=\tiny},
	every axis label/.append style = {font=\scriptsize}
}
\usepgfplotslibrary{fillbetween}
\pgfplotsset{compat=1.10}

\newtcolorbox[auto counter]{formule}{%
	skin=standard jigsaw,
	fonttitle=\bfseries\upshape,
	arc=2mm, opacityback=0
}

\definecolor{nocolor}{HTML}{282727}